%% SCRAP: architecture/PHYSICS_CONE_OF_INFLUENCE_DESIGN
%% SOURCE: docs/working/architecture/PHYSICS_CONE_OF_INFLUENCE_DESIGN.adoc
%% STATUS: WORKING
%% FITS: dev-guide/app-physics
%% EDITORIAL: lifted — prose rewritten to press voice
%% PATENT: This scrap describes adaptive feedback-control mechanisms. Multiple
%% passages below read as claim-like ("the system transforms async messaging
%% into deterministic ordering", "a closed feedback loop that narrows the cone").
%% Flagged inline with %% PATENT: markers. No claims are drafted here; forward to
%% patent counsel before any external release.

\section{The Cone of Influence}

The cone of influence is a feedback-control architecture for the physics engine.
It is named for a geometric metaphor used purely as a modeling device: the region
of possible execution futures is treated as a cone that narrows over time. As the
runtime accumulates observations, asynchronous and unpredictable publish/subscribe
messaging is progressively constrained into deterministic, ordered information flow.

%% PATENT: claim-like phrasing follows.
The driving idea is a closed loop. Observables from the early phase inform
statistical models. Those models adjust runtime knobs. The knobs change behavior.
The changed behavior generates new observables. Iterated, the loop moves the
system from high-uncertainty asynchronous communication toward low-uncertainty
deterministic ordering.

\subsection{The Model}

The cone occupies a conceptual three-axis space: a time axis (past to future), a
scope axis (narrow word-level interactions to wide global scope), and a
determinism axis (asynchronous fire-and-forget, through semi-deterministic
retry-with-hints, to deterministic causal ordering). The cone narrows along three
dimensions simultaneously:

\begin{itemize}
  \item \textbf{Temporal} --- past observations constrain present behavior.
  \item \textbf{Spatial} --- an initially global scope resolves to specific word
        interactions.
  \item \textbf{Modal} --- asynchronous uncertainty resolves to deterministic
        guarantees.
\end{itemize}

\subsection{Three Operational Zones}

\paragraph{Zone 1 --- Wide and Async (Discovery).} During the first $N$
observations the system bootstraps. Many execution paths remain possible, entropy
variance is high, latency is unpredictable, and scope is global. Publish events
fire-and-forget into a ring buffer with no ordering guarantees. The goal is simply
to discover which topics words publish to. Overhead is minimal, on the order of
200~ns. Observables available: entropy changes, temperature jumps, host snapshots,
error rates.

\paragraph{Zone 2 --- Narrowing and Semi-Deterministic (Learning).} After $M$
observations, entropy variance stabilizes, temperature reaches equilibrium, and
the latency distribution narrows. Scope contracts to module level. Publish events
now carry causal ordering hints: each message is tagged with a Lamport clock and a
likely-destination set, and related messages are batched. The goal is to learn the
dependency graph and subscriber patterns. Overhead rises to roughly 250~ns.
New observables: entropy slope, temperature stability, latency percentiles (p95,
p99), and the causal dependency graph.

\paragraph{Zone 3 --- Narrow and Deterministic (Enforcement).} In steady state,
entropy plateaus, temperature is bounded by throttling, and latency carries
guarantees against a service-level objective (SLO). Scope is specific: a word sends
only to known targets. Publish events are ordered deterministically --- buffered
and sorted by causal order number --- and a batch deadline of at most $T$~ms is
enforced. The p99 latency SLO is guaranteed, with backpressure applied on
violation, and a subscriber whitelist eliminates surprises. Overhead settles at
roughly 300~ns. The full observable set adds latency variance, thermal slope, and
execution-pattern periodicity.

\subsection{Metrics Inventory}

The engine already captures a substantial baseline. Per-word metrics include
\texttt{DictEntry.entropy} (raw execution count), and within
\texttt{DictEntry.physics}: \texttt{temperature\_q8} (EMA hotness),
\texttt{last\_active\_ns} (recency), \texttt{mass\_bytes} (footprint),
\texttt{avg\_latency\_ns} (rolling average), and \texttt{state\_flags}. Two fields,
\texttt{acl\_hint} and \texttt{pubsub\_mask}, are reserved stubs.

Host-level snapshots feed the analytics heap from real POSIX and procfs sources:

\begin{center}
\begin{tabular}{lll}
\toprule
Signal & Source & Status \\
\midrule
\texttt{monotonic\_time\_ns} & \texttt{clock\_gettime(CLOCK\_MONOTONIC)} & Real \\
\texttt{cpu\_count} & \texttt{sysconf(\_SC\_NPROCESSORS\_ONLN)} & Real \\
\texttt{scheduler\_policy} & \texttt{sched\_getscheduler(0)} & Real \\
\texttt{scheduler\_quantum\_ns} & \texttt{sched\_rr\_get\_interval()} & Real \\
\texttt{load\_avg\_milli} & \texttt{getloadavg()} & Real \\
\texttt{psi\_cpu/io/mem\_*\_milli} & \texttt{/proc/pressure/*} & Real (4.20+) \\
\texttt{cpu\_total/idle\_jiffies} & \texttt{/proc/stat} & Real \\
\texttt{cgroup\_cpu\_usage\_us} & \texttt{/sys/fs/cgroup/cpu.stat} & Real (cgroup v2) \\
\texttt{cgroup\_memory\_current} & \texttt{/sys/fs/cgroup/memory.current} & Real (cgroup v2) \\
\bottomrule
\end{tabular}
\end{center}

Profiler data (call frequency, total time, latency distribution, memory
read/write patterns) and test-infrastructure outcomes round out the current set.

Several signals are not yet captured but are reachable: stack and return-stack
depth at execution, dictionary lookup latency, block I/O queue depth, the data-flow
graph of word-to-word calls, and the error rate. Useful ML features remain on the
horizon: execution-pattern history, latency percentiles, entropy and thermal
slopes, instruction-count estimates, and cache-miss rates where performance
counters exist. Some quantities are out of reach by definition --- future execution
patterns require prediction, optimal scheduling is NP-hard, and ground-truth
causality runs into observer effects.

\subsection{The Control Layer: Runtime Knobs}

Per-word knobs include execution priority (tuned by temperature), cache affinity
(by entropy slope), temperature target, entropy half-life (cooling rate), latency
limit (the SLO), stack-depth limit (recursion guard), and ACL level. System-level
knobs include sampling cadence, analytics-heap size, profiler mode, pub/sub batch
size and timeout, temperature smoothing alpha, and entropy decay rate. Each knob is
driven by a specific tuning signal --- for example, the latency limit responds to
p99 measurements, and the sampling cadence responds to latency variance.

\subsection{The Pub/Sub Gateway: Three Phases}

The gateway is where words publish observations and where the control system acts.
Its behavior tracks the three zones.

In \textbf{Discovery}, a word calls
\texttt{physics\_analytics\_publish\_event(topic, payload, size)} and the event is
hash-inserted into the ring buffer; consumers poll when ready, with no ordering.

\begin{lstlisting}[language=C]
// Zone 1: fire-and-forget
physics_analytics_publish_event(PUBSUB_TOPIC_X, payload, size);
\end{lstlisting}

In \textbf{Learning}, the gateway identifies topic dependencies, learns subscriber
sets, stamps each message with a Lamport clock, and tags it with a subscriber hint.

\begin{lstlisting}[language=C]
// Zone 2: learning phase
message.causal_order_num = ++g_causal_clock;
message.subscriber_hint  = learned_targets;  // "probably {WordY, WordZ}"
physics_analytics_publish_event(PUBSUB_TOPIC_X, message, sizeof(message));
\end{lstlisting}

In \textbf{Enforcement}, causal ordering is enforced by buffering and sorting on
the order number; a batch deadline flushes pending messages; the p99 SLO is
verified; delivery is restricted to a whitelist; and a full queue triggers
backpressure on the publisher.

\begin{lstlisting}[language=C]
// Zone 3: enforcement phase
message.causal_order_num   = ++g_causal_clock;
message.subscriber_targets = learned_whitelist;   // exact set, no surprises
message.deadline_ns        = now_ns + SLO_LIMIT_NS;
\end{lstlisting}

\subsection{The Feedback Loop}

%% PATENT: claim-like phrasing follows.
The control law maps observed state, through a detection rule, to a knob
adjustment. A very hot word with rising entropy lowers execution priority, shortens
the entropy half-life, tightens the stack limit, and raises the sampling cadence. A
p99 latency that persistently exceeds the SLO tightens the stack-depth limit, the
batch timeout, and the batch size. A reached entropy plateau --- low slope, low
variance --- raises cache affinity, the temperature target, and execution priority.
Excessive latency variance raises sampling cadence and profiler detail. An error-rate
spike restricts the ACL level and lowers priority. Scheduler noise from core
migration raises cadence and pins affinity.

A hot-word emergency illustrates the loop end to end: a word saturates at
\texttt{0xFFFF} with a climbing entropy slope; the controller lowers OS priority,
accelerates cooling, caps recursion, raises sampling, and isolates the word via the
ACL level; the next cycle checks whether the temperature stabilized, restoring
normal priority on success or escalating to throttling and governance alert on
failure; offline analysis later mines the event for better detection heuristics.

\subsection{ML and Statistical Feedback (Future)}

Once the three zones operate and history accumulates, the design anticipates four
further phases: predictive models for latency, hotness, and optimal batch size;
anomaly detection for execution-pattern shifts, phase changes, SLO violations, and
early thermal runaway; optimization of core affinity, batching strategy, priority
ordering, and profiler mode; and governance integration with machine-checked proofs
in Isabelle --- observable invariants, a cone-narrowing convergence theorem,
verified control loops, and proven SLO bounds.

\subsection{Open Design Questions}

The design records several unresolved choices, each with a working recommendation:

\begin{itemize}
  \item \textbf{Zone transition} --- time-based, confidence-based, or hybrid.
        Recommendation: hybrid (start time-based, add confidence).
  \item \textbf{Ordering strength} --- causal only, total, or causal-within-module
        and total-globally. Recommendation: hybrid.
  \item \textbf{SLO scope} --- global, per-word, or governance-derived per category.
        Recommendation: derived.
  \item \textbf{Loop style} --- open-loop, closed-loop, or hybrid. Recommendation:
        hybrid, with closed-loop reserved for critical knobs.
  \item \textbf{ML training} --- in-system, offline, or hybrid. Recommendation:
        hybrid, with online anomaly response and offline refinement.
\end{itemize}

\subsection{Glossary}

\begin{itemize}
  \item \textbf{Causal ordering} --- messages ordered by dependency ($A \to B$ if
        $A$ must precede $B$).
  \item \textbf{Lamport clock} --- a scalar clock that respects causality.
  \item \textbf{Cone of influence} --- the region of possible execution futures,
        narrowing as confidence grows.
  \item \textbf{Zone} --- an operating region (async discovery, semi-deterministic
        learning, deterministic enforcement).
  \item \textbf{Knob} --- a runtime-adjustable parameter affecting behavior.
  \item \textbf{SLO} --- service-level objective (a latency guarantee).
  \item \textbf{Backpressure} --- flow control that slows publishers when a queue
        fills.
  \item \textbf{Entropy} --- execution count per word, an implicit hotness measure.
  \item \textbf{Temperature} --- the exponential moving average of entropy, an
        explicit hotness measure.
  \item \textbf{Half-life} --- the time for temperature to decay to half its peak.
\end{itemize}
